\chapter{Requirements and Design Decisions}
\label{ch:requirements}

This chapter recalls the functional requirements, shows that each of them is
implemented and tested, and explains the main design decisions they induced.

\section{Functional Requirements}
The core requirements (R1--R10) describe a checkout that treats purchases as a
set of bought items in mandatory categories (R1, R2, R2b), where each item is
priced and the price may depend on the quantity (R3), and payment is made by
bank card (R4). The system must offer \emph{extensible} discount plans --- at
least \emph{prime} (20\% off when the total reaches \euro{}50) and \emph{platinum}
(30\% off always) --- and let the owner introduce new ones (R5, R5b). It must
apply \emph{category-level pricing policies} that are likewise extensible (R6,
R6b). Customers may request home delivery (R7), charged according to weight and
distance, with refusal above 50\,kg (R8) and plan-based reductions (R8b).
Finally, the system maintains a live inventory that raises low-stock alerts
using the Observer pattern (R9) and supports two-hour delivery slots with
anti-overbooking and dynamic pricing (R10). Part~2 of the assignment adds a
command-line user interface (CLUI) and mandatory tests.

\section{Requirements Traceability}
\label{sec:traceability}
Table~\ref{tab:traceability} shows that every element required by the
instructions is implemented and tested. It maps each requirement to the classes
that realise it and to the JUnit tests (or scenario) that validate it.

\begin{table}[H]
\centering
\small
\begin{tabularx}{\linewidth}{p{1.6cm} Y Y}
\toprule
\textbf{Req.} & \textbf{Implemented in} & \textbf{Validated by} \\
\midrule
R1, R2, R2b & \code{Cart}, \code{CartEntry}, \code{Item}, \code{Category}; \code{Supermarket.setup} & \code{CliDispatcherTest}, \code{BillComputationTest} \\
R3 & \code{Item.unitPriceFor} (quantity-based price) & \code{ItemTest} \\
R4 & \code{POSDevice}, \code{TransactionAuthorisationSystem}, \code{BankCard} & \code{PaymentFlowTest}, \code{TransactionAuthorisationSystemTest} \\
R5 & \code{DiscountPlan} + Normal/Prime/Platinum & \code{DiscountPlanTest}, \code{BillComputationTest} \\
R5b & \code{DiscountPlanFactory}, \code{getAnnualFee} & \code{DiscountPlanTest}, \code{SupermarketTest} \\
R6, R6b & \code{PricingPolicy}, \code{CategoryDiscount} & \code{PricingPolicyTest} \\
R7 & \code{DeliveryRequest}, \code{Supermarket.requestDelivery} & \code{DeliveryCheckoutTest} \\
R8 & \code{StandardDeliveryCalculator} & \code{DeliveryCalculatorTest} \\
R8b & \code{DiscountPlan.applyDeliveryDiscount} & \code{DiscountPlanTest}, \code{DeliveryCheckoutTest} \\
R9 & \code{Inventory}, \code{StockObserver} (Observer) & \code{InventoryObserverTest} \\
R10 & \code{TimeSlot}, \code{DeliveryManager} & \code{DeliveryManagerTest}, \code{TimeSlotTest} \\
Customer ID (\S2.3) & \code{Customer.numericalId} & \code{SupermarketTest} \\
CLUI (Part 2) & \code{cli} package (login, setup, scanItem, pay, runTest, \ldots) & \code{CliDispatcherTest} \\
Test scenario (\S4) & \code{testScenario1.txt}, \code{my\_supermarket.ini} & \code{CliDispatcherTest} \\
\bottomrule
\end{tabularx}
\caption{Requirements traceability: every required feature is implemented and tested.}
\label{tab:traceability}
\end{table}

\section{How Requirements Shaped the Design}
Two words in the instructions drove most of the design: \emph{extensible} and
\emph{Observer}. The requirement to add new discount plans (R5b) and new pricing
policies (R6b) without modifying existing code led directly to the
\textbf{Strategy} pattern: discount plans and pricing policies are interfaces
with interchangeable implementations, so a new plan or policy is a new class
rather than an edit to the checkout. The same reasoning applies to the delivery
charging scheme (R8). The explicit demand for the \textbf{Observer} pattern in
R9 shaped the inventory module, where the stock count is decoupled from whoever
reacts to a shortage.

Three further decisions are worth noting. First, the bank-side protocol (R4) is
hidden behind a single \emph{authorise} operation, a \textbf{Facade} that keeps
the rest of the system independent of the banking details. Second, the code is
organised in layers (domain entities, services, orchestration, user interface)
so that the business logic does not depend on the way it is driven --- a property
that proved valuable when a second interface was later added. Third, all
monetary output is produced with a fixed locale (period decimals, UTF-8), so the
computed figures are identical on any machine that grades the project.

\section{The Bill Computation Pipeline}
\code{CashRegister.computeBill()} is the heart of the specified system. For the
items in the cart it applies, in order:
\begin{enumerate}[nosep]
  \item the \textbf{quantity-based unit price} (R3), through
  \code{Item.unitPriceFor(qty)};
  \item the \textbf{category pricing policy} (R6), e.g.\ a 10\% reduction on
  fruit and vegetables;
  \item the customer's \textbf{discount plan} (R5) on the resulting subtotal ---
  prime applies its 20\% only when that subtotal reaches \euro{}50, platinum
  always applies 30\%;
  \item the \textbf{delivery fee} (R8) when a delivery was requested, reduced by
  the plan (R8b: prime pays half, platinum free).
\end{enumerate}
The total is thus \emph{items after plan} plus \emph{plan-adjusted delivery}.
Subscribing to a plan also charges its annual fee immediately (normal
\euro{}0, prime \euro{}50, platinum \euro{}200), as required by \S2.3.
